\section{\paraphrase{inconsequential} Design choices} % MTODO: suggest another title using mparaphrase
\label{design-choices}
Several design choices in ILD's semantics are worth noting:
\begin{itemize}
\item The syntax shorthands for the \ildsf{quote} and \ildsf{macroexpand}
special forms (\cref{syntax}) desugar to the \ildsym{quote} and \ildsym{macroexpand}
symbols instead of the \ildsf{quote} and \ildsf{macroexpand} special form objects.
This restricts the parser output to simple S-expressions and allows the syntax to be
overridden by the program, but requires the respective keywords to be bound in the local
scope (also discussed in \cref{macro-function-duality}).
\item Bodies of \ildsetAbstr-terms (\cref{abstraction}, lambda-expressions)
carry one expression rather than
a list of expressions. This choice is consistent with ILD's pure-by-default approach.
Sequential execution of statements can instead, for example, be implemented as a \ild{do}-macro.
\item Capturing the local environment with a special form (\ildsf{free-vars}) instead
of making it available to every applicative (similar to fexpr-based languages \cite{kernel-report})
allows for a simpler interpreter and simpler (future) implementation of a partial evaluator.
\end{itemize}

